% GENERATED FILE. Edit canonical lessons, then run npm run generate:books.
% canonical-source-hash: fnv1a64:33d1f430400d7f45
% canonical-lessons: HI-C466-nisthur, HI-C466-krur, HI-C466-nek, HI-C466-sajjan, HI-C466-kathor

\chapter{Kind and Cruel}
\label{ch:hi-kind-and-cruel}
\hlchaptermodality{466}

\begin{chapteropening}
\textbf{By the end of this chapter:} \emph{I can say whether someone is kind or harsh.}

Complete the last lesson of chapter 466: I can say whether someone is kind or harsh.
\end{chapteropening}

\section[\texorpdfstring{\hi{निष्ठुर} (niṣṭhur)}{निष्ठुर (niṣṭhur)}]{\hi{निष्ठुर} (niṣṭhur) — heartless; unfeeling}

\label{lesson:HI-C466-nisthur}



\begin{quote}
Before the new one: say the Hindi for a sigh, then the Hindi for wailing.
\end{quote}

\subsection*{You'll want to know: \hi{निष्ठुर}}
\textbf{\hi{निष्ठुर}} — \emph{niṣṭhur} — \textquotedblleft{}heartless; unfeeling\textquotedblright{}.

\begin{quote}
\hi{बच्चा} \hi{रो} \hi{रहा} \hi{था}, \hi{लेकिन} \hi{वह} \hi{निष्ठुर} \hi{आदमी} \hi{नहीं} \hi{रुका।}
\end{quote}

\begin{quote}
\emph{baccā ro rahā thā, lekin vah niṣṭhur ādmī nahīṃ rukā.}
\end{quote}

\begin{quote}
The child was crying, but that heartless man didn't stop.
\end{quote}

\begin{cousinweb}
\textbf{\hi{निष्ठुर}} (\emph{niṣṭhur}) is a tatsama, borrowed from Sanskrit \emph{niṣṭhura}, \textquotedblleft{}hard, rough, harsh, cruel\textquotedblright{}, said especially of words (Monier-Williams).
\end{cousinweb}

\begin{grammarlens}[title={What you've built}]
One more word for describing what you lived and felt.
\end{grammarlens}

\subsection*{Your turn}
\begin{itemize}
\raggedright
  \item \emph{Say it:} \emph{niṣṭhur}
  \item \emph{Say it:} \emph{niṣṭhur}, once more
  \item \emph{Say it:} \emph{baccā ro rahā thā, lekin vah niṣṭhur ādmī nahīṃ rukā.} — The child was crying, but that heartless man didn't stop.
  \item \emph{Recall:} say the Hindi for a sigh, then the Hindi for wailing, then say \emph{niṣṭhur} again
\end{itemize}

\begin{tcolorbox}[breakable,colback=teal!4,colframe=teal!35!black,title={Before you move on}]
What does \hi{निष्ठुर} mean? (\textquotedblleft{}Heartless; unfeeling\textquotedblright{}.) Say it once more.
\end{tcolorbox}
\section[\texorpdfstring{\hi{क्रूर} (krūr)}{क्रूर (krūr)}]{\hi{क्रूर} (krūr) — cruel}

\label{lesson:HI-C466-krur}



\begin{quote}
Before the new one: say the Hindi for wailing, then the Hindi for heartless.
\end{quote}

\subsection*{You'll want to know: \hi{क्रूर}}
\textbf{\hi{क्रूर}} — \emph{krūr} — \textquotedblleft{}cruel\textquotedblright{}.

\begin{quote}
\hi{कुत्ते} \hi{को} \hi{मारना} \hi{एक} \hi{क्रूर} \hi{काम} \hi{है।}
\end{quote}

\begin{quote}
\emph{kutte ko mārnā ek krūr kām hai.}
\end{quote}

\begin{quote}
Beating a dog is a cruel thing to do.
\end{quote}

\begin{cousinweb}
\textbf{\hi{क्रूर}} (\emph{krūr}) is a tatsama, borrowed from Sanskrit \emph{krūra}, first \textquotedblleft{}wounded, sore\textquotedblright{} and then \textquotedblleft{}cruel, harsh\textquotedblright{} (Monier-Williams).
\end{cousinweb}

\begin{grammarlens}[title={What you've built}]
One more word for describing what you lived and felt.
\end{grammarlens}

\subsection*{Your turn}
\begin{itemize}
\raggedright
  \item \emph{Say it:} \emph{krūr}
  \item \emph{Say it:} \emph{krūr}, once more
  \item \emph{Say it:} \emph{kutte ko mārnā ek krūr kām hai.} — Beating a dog is a cruel thing to do.
  \item \emph{Recall:} say the Hindi for wailing, then the Hindi for heartless, then say \emph{krūr} again
\end{itemize}

\begin{tcolorbox}[breakable,colback=teal!4,colframe=teal!35!black,title={Before you move on}]
What does \hi{क्रूर} mean? (\textquotedblleft{}Cruel\textquotedblright{}.) Say it once more.
\end{tcolorbox}
\section[\texorpdfstring{\hi{नेक} (nek)}{नेक (nek)}]{\hi{नेक} (nek) — good-hearted; upright}

\label{lesson:HI-C466-nek}



\begin{quote}
Before the new one: say the Hindi for heartless, then the Hindi for cruel.
\end{quote}

\subsection*{You'll want to know: \hi{नेक}}
\textbf{\hi{नेक}} — \emph{nek} — \textquotedblleft{}good-hearted; upright\textquotedblright{}.

\begin{quote}
\hi{अरुण} \hi{के} \hi{दादा} \hi{बहुत} \hi{नेक} \hi{इंसान} \hi{थे।}
\end{quote}

\begin{quote}
\emph{Aruṇ ke dādā bahut nek insān the.}
\end{quote}

\begin{quote}
Arun's grandfather was a very good-hearted man.
\end{quote}

\begin{cousinweb}
\textbf{\hi{नेक}} (\emph{nek}) is usually traced to Persian \emph{nēk}, \textquotedblleft{}good\textquotedblright{}.
\end{cousinweb}

\begin{grammarlens}[title={What you've built}]
One more word for describing what you lived and felt.
\end{grammarlens}

\subsection*{Your turn}
\begin{itemize}
\raggedright
  \item \emph{Say it:} \emph{nek}
  \item \emph{Say it:} \emph{nek}, once more
  \item \emph{Say it:} \emph{Aruṇ ke dādā bahut nek insān the.} — Arun's grandfather was a very good-hearted man.
  \item \emph{Recall:} say the Hindi for heartless, then the Hindi for cruel, then say \emph{nek} again
\end{itemize}

\begin{tcolorbox}[breakable,colback=teal!4,colframe=teal!35!black,title={Before you move on}]
What does \hi{नेक} mean? (\textquotedblleft{}Good-hearted; upright\textquotedblright{}.) Say it once more.
\end{tcolorbox}
\section[\texorpdfstring{\hi{सज्जन} (sajjan)}{सज्जन (sajjan)}]{\hi{सज्जन} (sajjan) — a gentleman; a decent person}

\label{lesson:HI-C466-sajjan}



\begin{quote}
Before the new one: say the Hindi for cruel, then the Hindi for good-hearted.
\end{quote}

\subsection*{You'll want to know: \hi{सज्जन}}
\textbf{\hi{सज्जन}} — \emph{sajjan} — \textquotedblleft{}a gentleman; a decent person\textquotedblright{}.

\begin{quote}
\hi{दरवाज़े} \hi{पर} \hi{एक} \hi{सज्जन} \hi{आपसे} \hi{मिलना} \hi{चाहते} \hi{हैं।}
\end{quote}

\begin{quote}
\emph{darvāze par ek sajjan āpse milnā cāhte haiṁ.}
\end{quote}

\begin{quote}
There's a gentleman at the door who would like to see you.
\end{quote}

Masculine.

\begin{cousinweb}
\textbf{\hi{सज्जन}} (\emph{sajjan}) is a tatsama, borrowed from Sanskrit \emph{saj-jana}, \textquotedblleft{}a good or virtuous man\textquotedblright{}, built on \emph{sat}, \textquotedblleft{}good\textquotedblright{}, and \emph{jana}, \textquotedblleft{}person\textquotedblright{} (Monier-Williams). It is the polite word for a man one does not know.
\end{cousinweb}

\begin{grammarlens}[title={What you've built}]
One more word for describing what you lived and felt.
\end{grammarlens}

\subsection*{Your turn}
\begin{itemize}
\raggedright
  \item \emph{Say it:} \emph{sajjan}
  \item \emph{Say it:} \emph{sajjan}, once more
  \item \emph{Say it:} \emph{darvāze par ek sajjan āpse milnā cāhte haiṁ.} — There's a gentleman at the door who would like to see you.
  \item \emph{Recall:} say the Hindi for cruel, then the Hindi for good-hearted, then say \emph{sajjan} again
\end{itemize}

\begin{tcolorbox}[breakable,colback=teal!4,colframe=teal!35!black,title={Before you move on}]
What does \hi{सज्जन} mean? (\textquotedblleft{}A gentleman; a decent person\textquotedblright{}.) Say it once more.
\end{tcolorbox}
\section[\texorpdfstring{\hi{कठोर} (kaṭhor)}{कठोर (kaṭhor)}]{\hi{कठोर} (kaṭhor) — harsh; stern}

\label{lesson:HI-C466-kathor}



\begin{quote}
Before the new one: say the Hindi for good-hearted, then the Hindi for a gentleman.
\end{quote}

\subsection*{You'll want to know: \hi{कठोर}}
\textbf{\hi{कठोर}} — \emph{kaṭhor} — \textquotedblleft{}harsh; stern\textquotedblright{}.

\begin{quote}
\hi{पिता} \hi{की} \hi{बात} \hi{कठोर} \hi{थी}, \hi{लेकिन} \hi{सही} \hi{थी।}
\end{quote}

\begin{quote}
\emph{pitā kī bāt kaṭhor thī, lekin sahī thī.}
\end{quote}

\begin{quote}
Father's words were harsh, but they were right.
\end{quote}

\begin{cousinweb}
\textbf{\hi{कठोर}} (\emph{kaṭhor}) is a tatsama, borrowed from Sanskrit \emph{kaṭhora}, \textquotedblleft{}hard, solid, stiff\textquotedblright{} (Monier-Williams). Hindi uses it for hard words and hard hearts.
\end{cousinweb}

\begin{grammarlens}[title={What you've built}]
One more word for describing what you lived and felt.
\end{grammarlens}

\subsection*{Your turn}
\begin{itemize}
\raggedright
  \item \emph{Say it:} \emph{kaṭhor}
  \item \emph{Say it:} \emph{kaṭhor}, once more
  \item \emph{Say it:} \emph{pitā kī bāt kaṭhor thī, lekin sahī thī.} — Father's words were harsh, but they were right.
  \item \emph{Recall:} say the Hindi for good-hearted, then the Hindi for a gentleman, then say \emph{kaṭhor} again
\end{itemize}

\begin{tcolorbox}[breakable,colback=teal!4,colframe=teal!35!black,title={Before you move on}]
What does \hi{कठोर} mean? (\textquotedblleft{}Harsh; stern\textquotedblright{}.) Say it once more.
\end{tcolorbox}
